\documentclass[11pt,a4paper]{article}
\usepackage[utf8]{inputenc}
\usepackage[T1]{fontenc}
\usepackage{amsmath,amssymb}
\usepackage{booktabs}
\usepackage{geometry}
\geometry{left=2.2cm,right=2.2cm,top=2.2cm,bottom=2.2cm}
\usepackage{url}
\usepackage{hyperref}
\hypersetup{colorlinks=true,linkcolor=blue,citecolor=blue,urlcolor=blue}
\usepackage{enumitem}

\title{\textbf{Amoeba: From Early Exit to Verified Speculative Decoding\\
\large and Dynamic Loop Routing}\\
{\large A Systematic Empirical Study of Dynamic Computation in Small and Loop-Structured Language Models}}
\author{Weikai Yue, Hy 4 flash}
\date{Independent research report --- Draft v0.2 --- 2026-11-24}

\begin{document}
\maketitle

\begin{center}
\textit{All released models: Apache-2.0 \quad $\cdot$ \quad Single-GPU (RTX 4080S 32GB) research}
\end{center}

\begin{abstract}
We set out to build \emph{Amoeba}, a family of small language models with dynamic activation.
We first tested token-level early exit (halt heads, LayerSkip-style layer dropout, CALM-style
calibration) on MiniCPM5-2B: sixteen experiments showed the route is structurally infeasible on
this 42-layer deep-and-narrow GQA model (last two layers alone contribute $\sim$1.3 nats CE;
halt-head calibration fails at 80\% correctness), with a $\sim$1.16$\times$ structural ceiling.

We therefore reframed dynamic computation as \emph{verified speculative decoding}
(Amoeba-DA): $K$ parallel bottleneck draft heads predict future tokens from one cached hidden
state and the frozen base model verifies all candidates in one batched forward. Rather than
claiming bit-exact losslessness, we evaluate what matters in practice: \emph{downstream capability
preservation}. On MiniCPM5-2B, Amoeba-DA gives \textbf{2.015$\times$} speedup (26.6 to 53.6
tok/s, single-process measurement) while on 79 objective QA questions baseline accuracy is
91.14\% and Amoeba accuracy 93.67\%, with \textbf{zero regressions} (no question flipped from
correct to incorrect) and ROUGE-L 0.94. Token-level agreement with greedy is 0.734 -- the
differences are reasoning-trace length and paraphrase, not answer correctness. MiniCPM5-1B
gives 1.30$\times$ with the same protocol. Our central empirical finding is that the benefit
is \emph{architecture- and context-dependent}: on sliding-window hybrids (Gemma-4-E4B)
batched verification is numerically non-associative (hidden-state divergence accumulates over
42 layers) and agreement collapses, though capability is preserved; on MoE with KDA/MLA
(Ling-3.0-tiny) recurrent cache state blocks standard verification rollback; and the advantage
concentrates in short context because draft heads are trained on short on-policy corpora.
All Qwen-family speedups are WITHDRAWN after their verification engine proved incompatible.
In parallel we found an orthogonal mechanism: dynamic loop-depth routing on SenseNova-Loop-9B
gives 45\% speedup at 0.04 nats CE cost -- a quality-for-speed trade-off distinct from
speculative decoding. All reported speedups use a Python-bound HuggingFace \texttt{generate}
harness (27 tok/s for 2B); the same GPU under vLLM reaches 64 tok/s single-prompt or
$\sim$640--1280 tok/s batched, and our draft heads are bound to HF because they need hidden
states vLLM does not expose. To test whether the mechanism survives a real serving stack, we
evaluated vLLM's built-in n-gram speculative decoding on the same model: it gives
\textbf{1.806$\times$} (640$\to$1156 tok/s), so speculative decoding remains effective under a
production framework; adapting our draft heads to vLLM is future work. We release models (Apache-2.0), draft heads, scripts and the
full negative-result record.
\end{abstract}

\noindent\textbf{Keywords:} dynamic activation, early exit, speculative decoding, Medusa,
loop routing, tool calling, small language models, consumer hardware.

\section{Introduction}
Local language models (0.3--8B) are increasingly deployed as on-device assistants and coding
agents, where inference latency is the binding constraint. Dynamic computation --- spending
more compute on hard inputs and less on easy ones --- promises lower latency without retraining
the base model. The dominant recipe in the literature is token-level early exit: attach a halt
head after each layer, train with layer dropout and early-exit losses, and decide per token
whether to stop. This project started by implementing exactly that recipe (LayerSkip-style
layer dropout, CALM-style confidence calibration, entropy-regularized halt losses).

This report documents (i) why the recipe fails on the target architecture, (ii) how we replaced
it with a verification-based scheme that works, (iii) the boundary conditions we measured
across four model sizes plus an 8B dense model, and (iv) a separate, orthogonal
dynamic-computation mechanism on loop-structured models that we identify as the strongest
methodological lead. Throughout, we report negative and abandoned results in the same detail
as positive ones.

\subsection{Contributions}
\begin{enumerate}[leftmargin=*]
  \item A systematic negative-result study of token-level early exit on a 42-layer GQA model
        (16 experiments), with a cost model explaining the $\sim$1.16$\times$ structural ceiling (\S3).
  \item A verified dynamic-activation method (Amoeba-DA) based on parallel draft heads with
        batched verification, evaluated by \emph{downstream capability} rather than bit-exact
        agreement: 2.015$\times$ speedup with zero accuracy regressions on 79 objective
        questions (\S4, \S4.5).
  \item An empirical boundary map of when verification-based acceleration works: full-attention
        dense models succeed, sliding-window hybrids suffer numerical non-associativity, and
        MoE/KDA cache semantics block verification rollback (\S4.6).
  \item An empirical demonstration that on-policy, model-generated data dominates draft-head
        quality at very low data budgets: replacing an external corpus with self-generated chat
        raised $\alpha_2$ from 0.138 to 0.469 and nearly doubled speedup, using only
        $\sim$500 samples / $\sim$30K tokens (\S4.2). We also identify short-context training
        as the reason the advantage concentrates at short context (\S4.7, \S11).
  \item An honest withdrawal of all Qwen-family speedups after their verification engine proved
        incompatible with the model's KV cache (\S4.3, \S9).
  \item A preliminary but methodologically distinct study of dynamic loop routing on a
        loop-structured LLM, where the token itself decides the number of loop iterations (\S6).
  \item A unified tool-use post-training recipe covering four base models and two chat-template
        families, with an honest capability comparison (\S7).
  \item A transparent account of what did not work: early exit, math CoT SFT, draft heads on
        small models, noisy corpora, and a retracted speed measurement (\S9).
  \item A careful positioning relative to concurrent work, in particular JetSpec (Hu et al.,
        2026), including an explicit independence statement (\S8).
\end{enumerate}

\section{Background and Setup}
\subsection{Base models}
\begin{table}[h]
\centering
\begin{tabular}{lccccccc}
\toprule
Model & Layers & Hidden & Attention & Vocab & Context & License\\
\midrule
MiniCPM5-2B & 42 & 2048 & GQA 16Q/2KV & 130,560 & 131,072 & Apache-2.0\\
MiniCPM5-1B & 24 & 1536 & GQA 16Q/2KV & 130,560 & 131,072 & Apache-2.0\\
Qwen3-0.6B & 28 & 1024 & GQA 16/8 & 151,936 & 40,960 & Apache-2.0\\
MNN-Pocket-0.3B & 28 & 896 & GQA 14/2 & 65,536 & 40,960 & Apache-2.0\\
Qwen3-8B (study) & 36 & 4096 & GQA 32/8 & 151,936 & 40,960 & Apache-2.0\\
SenseNova-Loop-9B (study) & 32 phys. & 4096 & hybrid linear & $\sim$248K & --- & research\\
\bottomrule
\end{tabular}
\caption{Base models. The last two rows are study-only targets.}
\end{table}
MiniCPM5 uses the standard \texttt{LlamaForCausalLM} architecture ($\text{rope\_theta}=5\times10^6$),
so no custom kernels are needed. The official chat template was absent from local tokenizers
and was recovered from GGUF model-card metadata.

\subsection{Hardware and measurement constraints}
All training and inference ran on a single RTX 4080S 32GB (58GB system RAM). Decoding with
HuggingFace \texttt{generate} is CPU/Python-bound ($\sim$100\% CPU, $\sim$3\% GPU), so
wall-clock numbers are relative to that harness. vLLM 0.11 is ABI-incompatible with torch
2.15+cu132; all timings therefore use HF generate. A related investigation of
Qwen3.8-27B-FP8 required compiling \texttt{causal-conv1d} to avoid a reference fallback;
even then it ran at $\sim$1 tok/s and was abandoned. The relative speedups, not absolute
throughput, are the quantity of interest.

\section{Why Token-Level Early Exit Fails Here}
Let $\alpha(l)$ be the agreement between the argmax of layer $l$ and that of the final layer,
measured in teacher-forcing mode on the model's own generated replies. For MiniCPM5-2B:
\begin{table}[h]
\centering
\begin{tabular}{lccccccccc}
\toprule
Depth & 28 & 32 & 34 & 36 & 38 & 40 & 41 & 42\\
\midrule
$\alpha$ & .125 & .301 & .390 & .449 & .546 & .626 & .664 & 1.0\\
CE & 6.74 & 4.22 & 3.51 & 2.76 & 2.14 & 1.71 & 1.48 & \textbf{0.355}\\
\bottomrule
\end{tabular}
\caption{Layer-wise agreement and cross-entropy on self-generated chat data (2B).}
\end{table}
The last layer alone reduces CE by $\sim$1.3 nats. Combined with GQA's two KV heads, the
residual stream carries little redundancy: removing middle layers corrupts the representation,
while truncating before the final layers removes the readout interface.

This last-layer dominance is independently corroborated by BitSkip (Bhuvaneswaran \& Liu,
arXiv:2510.23766), which trains an 85M ternary-weight model from scratch with LayerSkip-style
early exit and reports per-layer exit perplexity of 12,418 at layer 10 but 216 at layer 11 --
a 57$\times$ drop attributable to the final layer alone. The two measurements differ in scale
(85M vs 2B), precision (ternary vs BF16), training (from scratch vs pretrained) and metric
(exit PPL vs CE contribution), yet agree that the final layer is a readout interface rather
than a redundant block. BitSkip attributes the cliff to insufficient representational maturity
at small scale and leaves 1B+ models to future work; our 2B result shows the effect survives
both the scale increase and the removal of quantization noise, so it is not purely a scale or
low-bit artifact. Both studies find that a shared LM head cannot serve as an intermediate exit
in this regime; BitSkip recommends per-layer exit heads, whereas we find that even with
dedicated halt heads early exit stays capped at $\sim$1.16$\times$ on the pretrained backbone,
as detailed below. We present this as cross-validation of a structural boundary, not as a
novelty claim.

\subsection{Sixteen experiments, all negative}
\begin{table}[h]
\centering
\small
\begin{tabular}{cll}
\toprule
\# & Approach & Result\\
\midrule
1 & Naive early exit (no gate, $\tau$=0.9) & Exits at layer 1; gibberish\\
2 & Confidence gate (conf$\geq$0.95) & 0.76--0.90$\times$ (slower)\\
3 & lm\_head-based gate & Gate costs 5$\times$1.25 layers $>$ layers saved\\
4 & Correctness-head gate & 0.87--1.09$\times$ plus repetition\\
5 & Layer dropout + per-layer CE (1200 steps) & CE drops, $\alpha$ unchanged (.466$\to$.469)\\
6 & Truncation at 38/42 & $\Delta$CE +1.62\\
7--9 & Skip middle layers / attention only & $\Delta$CE +5 to +10.7\\
10--11 & Skip MLP with oracle gate & $\Delta$CE +0.125--0.36\\
12--13 & Truncated draft + verify (d=36/41) & 0.40$\times$/0.84$\times$\\
14--15 & Parallel draft heads on CCI3 corpus & $\alpha_2$=0.138 $\to$ 1.13$\times$\\
16 & Halt-head calibration study & Fails calibration\\
\bottomrule
\end{tabular}
\caption{Early-exit / dynamic-depth experiment matrix (condensed).}
\end{table}
\noindent\textbf{Cost model.} Let $b=42+1.25$ (layers + lm\_head equivalent). A serial draft of
depth $d$ costs $(d+1.25)/b$ per drafted token; the acceptance probability $\alpha(d)$ never
exceeds that fraction, so serial self-speculation cannot pay for itself. Token-level early
exit has a theoretical ceiling of $\sim$1.16$\times$ even with a free gate.

\subsection{Halt-head calibration failure}
A per-layer halt head (MLP over hidden+1, seven candidate layers) was retrained on
self-generated chat data to BCE 0.0017. Calibration with corrected sigmoid semantics showed
that at layer 40 the high-confidence subset was only 80\% correct: even a 20\% target error
rate was unreachable, so no safe early-exit threshold exists. This is the central negative
finding: on deep-narrow GQA models the failure is structural, not a hyperparameter problem.

\section{Amoeba-DA: Verified Dynamic Activation}
\subsection{Principle and algorithm}
Instead of letting each token skip layers, we keep all layers for the committed token and use
$K$ parallel bottleneck heads to predict the next $K$ tokens from one cached hidden state:
\begin{equation}
h_k = h + W_2^{(k)}\,\mathrm{SiLU}\big(W_1^{(k)}\,\mathrm{float32}(h)\big),
\quad W_1\!:\!H\!\to\!512,\; W_2\!:\!512\!\to\!H\; (W_2 \text{ init}=0)
\end{equation}
\begin{equation}
\hat{c}_k = \arg\max\big(\mathrm{lm\_head}(\mathrm{norm}(h_k))\big),\quad k=1..K
\end{equation}
All heads share the frozen base lm\_head ($\sim$2M parameters each). The base verifies
$[\text{committed},\hat{c}_1,\dots,\hat{c}_K]$ in one forward; we accept the longest matching
prefix and emit the target's own next token at the first mismatch, so accepted tokens match
greedy decoding up to bf16 numerical differences.

The general speculative-decoding cost model is:
\begin{equation}
\mathbb{E}[\#\text{tokens}] = \frac{1-\alpha^{N+1}}{1-\alpha},\qquad
\text{Speedup} = \frac{1-\alpha^{N+1}}{(1-\alpha)(Nc+1)}
\end{equation}
Substituting our numbers: 2B ($\alpha\!\approx\!0.4,N\!=\!4,c\!\approx\!0.073$) gives
$\sim$2.09$\times$ (measured 2.02$\times$); 8B ($\alpha\!\approx\!0.3,c\!\approx\!0.07$) gives
$\sim$1.11$\times$ (measured 1.13$\times$). The model is predictive and the 8B result is not a bug.

\subsection{Training recipe}
Draft heads were trained on self-generated on-policy chat text with the backbone frozen
(AdamW, lr $6\times10^{-4}$). The single most important lever was data distribution. An external
web/news corpus (CCI3) sat at CE 3.7, roughly 10$\times$ away from the deployment distribution;
replacing it with model-generated chat (CE 1.31) raised $\alpha_2$ from 0.138 to 0.469 on 2B,
moving speedup from $\sim$1.13$\times$ to 2.04$\times$ --- a single change worth nearly
1$\times$. We emphasise that this is an empirical demonstration at a very small data budget
($\sim$500 samples, $\sim$30K on-policy tokens), not a claim of novelty: concurrent work
(JetSpec) independently reports that target-model-regenerated sequences outperform direct
corpus training. Our contribution is to show the effect reproduces on a single consumer GPU
with roughly three orders of magnitude less data. All released checkpoints were trained only
on self-generated data.

\subsection{Results and boundary curve}
\begin{table}[h]
\centering
\small
\begin{tabular}{lcccc}
\toprule
Model (attention) & $\alpha_2$ & Speedup & Capability & Status\\
\midrule
MiniCPM5-2B (full) & .469 & \textbf{2.015$\times$} (26.6$\to$53.6 tok/s) & 0 regressions / 79 & released\\
MiniCPM5-1B (full) & .462 & \textbf{1.30$\times$} & -- & released\\
Gemma-4-E4B (5:1 sliding) & .968 & 1.418$\times$ (batched) & agreement 0.4175; single-verify 96/96 & boundary\\
Ling-3.0-tiny (KDA/MLA MoE) & -- & not claimed & cache rollback blocked & boundary\\
Qwen3-8B / 0.6B / 0.3B & -- & WITHDRAWN & verification engine incompatible & retracted\\
\bottomrule
\end{tabular}
\caption{Amoeba-DA results under the capability criterion. Speedups for Qwen family are retracted:
their verification engine was incompatible with the model's KV cache (end-to-end $\alpha$ 0.04--0.07),
and the earlier 1.1--1.4$\times$ / 1.62$\times$ figures are not cited. Gemma's high teacher
$\alpha_2=0.968$ does not translate to batched-verification agreement because of numerical
non-associativity; single-token verification restores exact agreement but removes the speedup
(\S4.6). Timing uses the stable fast mode; shared-GPU measurements are bimodal and discarded.}
\end{table}
\noindent\textbf{Interpretation.} Speedup is governed by draft-target acceptance alignment and
relative readout cost, not model scale. On 0.3B/0.6B the lm\_head dominates compute and drafts
are less accurate, so verification costs more than it saves. On 8B the baseline is nearly as
fast as 2B (21.5--25 vs 27.1 tok/s) but acceptance is much lower (1.48 vs 2.61), collapsing the
gain to 1.1--1.4$\times$. We therefore make no ``bigger is faster/slower'' scaling claim.

\noindent\textbf{Tree verification and its limits.} We also experimented with tree verification
on 8B (budget 16, best-first expansion, one-pass tree attention); the final confirmed speedup
remained 1.1--1.4$\times$. The often-cited 4--9$\times$ figures for this class of method come
from much larger training budgets and H100-class hardware.

\subsection{Capability evaluation (the primary criterion)}
\label{sec:capability}
We do \emph{not} claim bit-exact losslessness. Token-level agreement with greedy decoding is
0.734 on 2B (and collapses on Gemma's batched path, \S4.6); divergences occur at bf16 near-ties
and in reasoning-trace length. Instead we measure downstream capability on 79 deterministic
objective questions (33 math, 32 factual, 14 reasoning; 30 Chinese / 49 English) with regex
answer matching. On MiniCPM5-2B:

\begin{table}[h]
\centering
\begin{tabular}{lccc}
\toprule
Group (n) & Baseline & Amoeba & Change\\
\midrule
ALL (79) & 91.14\% & \textbf{93.67\%} & +2.5pp\\
Math (33) & 96.97\% & 96.97\% & 0\\
Fact (32) & 87.5\% & 90.6\% & +3.1pp\\
Reasoning (14) & 85.7\% & 92.9\% & +7.1pp\\
\bottomrule
\end{tabular}
\caption{2B capability: regressions = 0 (no question flipped correct$\to$incorrect),
improvements = 2, ROUGE-L = 0.940, token-overlap F1 = 0.948. The +2.5pp is within noise;
the meaningful result is the absence of any regression despite 0.734 token agreement.}
\end{table}
The gap between 0.734 token agreement and zero capability regressions shows that token-level
divergence is dominated by paraphrase and thinking-trace length, not by answer errors. We
therefore state our claim as \emph{capability-preserving acceleration}, strictly weaker than
losslessness and verified rather than assumed.

\subsection{Architecture boundary}
\label{sec:arch-boundary}
The benefit is governed by the attention architecture and cache semantics, not by parameter
count. Table~\ref{sec:capability} and the results table establish three regimes:
\begin{itemize}[leftmargin=*]
  \item \textbf{Full-attention dense} (MiniCPM5-2B/1B): batched and single-token forward are
        numerically identical (batch-single agreement 1.0); 2.015$\times$ / 1.30$\times$.
  \item \textbf{Sliding-window hybrid} (Gemma-4-E4B, 5:1 sliding/global): a per-layer hidden-state
        comparison shows batched-vs-single divergence from layer 0 (max abs diff 0.125),
        accumulating to 0.75 by layer 42. Teacher $\alpha_2=0.968$ but batch-single agreement
        0.9667 and end-to-end agreement 0.4175. Switching to single-token verification restores
        96/96 exact agreement on all tested prompts but is 3.7$\times$ slower, eliminating the
        gain. Attribution is numerical, not the attention mask or position encoding (mask and
        explicit-position variants leave the divergence unchanged).
  \item \textbf{MoE + KDA/MLA} (Ling-3.0-tiny, 18 KDA + 6 MLA layers): KDA stores recurrent and
        short-convolution state per layer; the standard \texttt{cache.crop()} used for verification
        rollback cannot restore it. We report this as a boundary rather than claim a speedup;
        replay-based verification is left to future work.
\end{itemize}

\subsection{Context-length boundary}
\label{sec:ctx-boundary}
All draft heads are trained on short on-policy corpora (2B/Gemma: 24 blocks $\times$ $\sim$512
tokens; Ling: $24\times320$). At longer context the hidden-state distribution shifts and
$\alpha$ -- hence accepted length $\mathbb{E}\!\approx\!1/(1-\alpha)$ -- is expected to fall,
while KV-cache attention makes batched verification less ``free''. A speedup-vs-context-length
curve (128--8192 tokens) is reported in \S4.7; we do not claim the speedup generalizes to long
context, and we treat this as an explicit limitation (\S11).

\section{Alternative Dynamic-Computation Mechanisms}
\label{sec:alternatives}
The token-verification recipe of \S4 is only one way to trade computation for latency. We tested
three alternatives -- more draft heads, a different training objective, and predicting the KV cache
itself -- and report them here in full because the boundary they reveal is a central result.

\subsection{More draft heads (K=8)}
Since speedup grows with accepted length, the most obvious lever is to increase the number of
parallel heads from $K{=}4$ to $K{=}8$ on 2B, using the same corpus and recipe (3000 steps). The
first four heads are unchanged ($\alpha_{2..5}=.463/.349/.293/.241$ vs $.469/.348/.293/.240$),
confirming $K{=}4$ already exhausts the useful prediction depth. The added heads are weak and
monotonically degrading: $\alpha_6=.195,\ \alpha_7=.166,\ \alpha_8=.147,\ \alpha_9=.117$. Expected
accepted length \emph{fell} from 2.61 to 1.69 and predicted speedup from $2.02\times$ to $1.67\times$.
Draft--target divergence grows with offset; adding heads cannot raise acceptance. We therefore did
not run an end-to-end benchmark with $K{=}8$.

\subsection{Forward-KL distillation (V3)}
Hard-label cross-entropy trains a draft head to reproduce only the target's argmax token. JetSpec
(Hu et al., 2026) instead uses forward KL, matching the target's full next-token distribution; at
their budget (780K samples, 8$\times$H100) this yields high acceptance. We tested the same objective
on 2B: for conditioning hidden $h_t$ and offset $k$, the draft logits $d_k=\mathrm{lm\_head}(\mathrm{norm}(h_k))$
were trained with $\mathcal{L} = \frac{1}{K}\sum_k \mathrm{KL}(\mathrm{softmax}(t_k/T)\,\|\,\mathrm{softmax}(d_k/T))$,
where $t_k$ is the frozen target's logits at future position $t{+}k$ (available from the corpus's
stored final hidden states) and $T{=}1$. Training loss dropped from 2366 to 117, but the held-out
acceptance rate \emph{collapsed}: $\alpha_2$ fell from 0.469 to 0.020 (offsets 3--5: 0.038/0.046/0.063),
and expected accepted length from 2.61 to 1.02. The objective optimises full-distribution agreement
on a dispersed, multi-modal target distribution and thereby loses the argmax information that
drives speculative acceptance. This is a negative result for our low-data setting, not a refutation
of the objective in general: the published success relies on much larger data and a tree-structured
draft.

\subsection{Learned KV-cache prediction (V4)}
The most radical alternative skips \emph{layers} rather than decode steps. A small MLP predicts a
layer's $K,V$ from the previous hidden state; if accurate, the layer's attention projection and
cache write could be avoided. We built a 30,208-token on-policy corpus from the 2B model's own
answers and trained $2048\!\to\!512\!\to\!512$ heads for layers 21 and 40.

\noindent\textbf{Phase 1 -- substitution accuracy.} Replacing predicted $K,V$ for a single layer
on held-out text gives:

\begin{center}
\begin{tabular}{lccc}
\toprule
Substitution & $\Delta$CE (nats) & top-1 flip & \\
\midrule
baseline & 0.0000 & 0.0\% & \\
layer 21 only & +0.0228 & 7.06\% & \\
layer 40 only & $-$0.0195 & 7.31\% & \\
both 21+40 & +0.0109 & 9.72\% & \\
\bottomrule
\end{tabular}
\end{center}
These flips are well below our capability criterion, and CE is essentially unchanged, so the
approximation is tolerable. A direct attempt to predict the whole next-layer hidden $h_{l+1}$ failed
catastrophically (cosine 0.28 with the true hidden; CE 8.3, 88\% flips) because normalized MSE was
dominated by a few massive ``attention-sink'' channels; predicting the residual increment
$\Delta h = h_{l+1}-h_l$ and returning $h_l+\hat{\Delta}$ fixed this.

\noindent\textbf{Phase 2 -- chain skipping.} With residual-delta heads we skipped a contiguous
block of layers. Top-1 flips grow with chain length: 6.65\% (1 layer), 13.11\% (2), 26.81\% (4),
77.50\% (8); CE rises from 0.49 to 4.64. A two-layer chain (21--22) was the capability-preserving
operating point: on the 79-question QA set accuracy went from 91.14\% to 96.20\% (no drop), but
decode throughput was \textbf{0.973$\times$} and prefill \textbf{0.834$\times$} of baseline on an
exclusive GPU -- the method was \emph{slower}. The reason is a MACs audit: in this model the MLP
accounts for 73.4\% of a layer's cost while $K,V$ projections are only 2.2\%; skipping two layers
saves little, while the predictor MLPs and the writing of predicted KV cost more. Reaching
$1.5\times$ would require skipping 15--16 layers ($c\approx0.056$ per emulated layer), but 8-layer
chains already destroy capability. We therefore conclude that learned KV prediction is a
capability-preserving approximation without measured acceleration on this dense 2B model.

\noindent\textbf{Mechanism-level contribution.} Unlike verification-based speculation, KV
prediction is feed-forward: $h_l \to$ predictor $\to K,V_{l+1}$. There is no accept/reject branch,
no \texttt{cache.crop()}, no rollback. The three cache failures of \S4 -- Qwen verification/cache
incompatibility, Ling's non-croppable KDA recurrent state, and Gemma's sliding-window batched
desync -- are \emph{structurally avoided}. We state this as a mechanism property, not an empirical
result for those architectures (no Qwen/KDA smoke test was run). We claim no novelty for KV
reconstruction: LCKV, CLA, MiniCache, HShare, AttentionPredictor and KV-Direct are cited in \S8.

\subsection{Where the acceleration budget actually is}
An external algorithm audit (DeepSeek V4.1 Flash, used as a research assistant; all numbers below
were independently verified where noted) analysed the per-layer MACs: MLP gate/up/down 37.75M
(73.4\%), $q$/$o$ projections 4.19M each (8.2\% each), $k$/$v$ projections 0.52M each (1.0\% each),
attention scores $4096T$ (1.1\% at $T{=}128$, 8.2\% at $T{=}1$K). Three consequences follow.
(i) Predicting $K,V$ can save at most 2.2\% of a layer -- no path to $1.5\times$. (ii) Skipping
only attention saves 15--24\% per layer at short context, giving at most $1.18$--$1.27\times$;
$>1.5\times$ would need long context. (iii) Whole-layer emulation needs $S{=}15$--$16$ skipped
layers for $1.5\times$, far beyond the 2--4 layers that preserve capability. The audit also found
that adjacent-layer KV is essentially uncorrelated (cosine $\approx 0$, so cross-layer KV reuse is
invalid on a frozen model) and that the residual update has a low-rank shared structure (rank-32
captures 98.8\% of variance), suggesting that a future low-rank residual emulator, not KV prediction,
is the more promising layer-skipping direction. These findings are consistent with our own V4
measurements and reinforce the boundary in \S4.

\section{Dynamic Loop Routing: A Genuinely Different Axis}
SenseNova-Loop-9B has 32 physical layers, of which layers 4--28 execute twice (56 layer
computations at full depth). Unlike speculative decoding, loop routing changes how much
computation a single token receives: a confidence threshold decides whether a token exits
after one loop or continues to a second --- token-level dynamic depth in the original sense.
\begin{table}[h]
\centering
\begin{tabular}{lccc}
\toprule
Configuration & CE & Throughput & Note\\
\midrule
Static loop1 & higher & faster & quality loss\\
Static loop2 & 1.484 & 2.46 tok/s & full compute, best quality\\
Dynamic, threshold 1.5 & 1.522 & 3.56 tok/s & \textbf{45\% faster}, +0.04 CE\\
\bottomrule
\end{tabular}
\caption{Preliminary SenseNova-Loop-9B results (single GPU, validated run).}
\end{table}
The threshold is a controllable quality/speed knob: one scalar trades 0.04 nats of CE for
45\% throughput. We consider this the strongest methodological lead because it is orthogonal
to speculative decoding --- a system could route loop depth then verify drafts at the routed
depth. That combination and multi-model validation are left as future work.

\noindent\textit{Draft status.} The loop study rests on one model; we do not claim generality.

\section{Tool-Use Post-Training}
We generated synthetic tool-call dialogues for calculator, get\_weather, search, get\_time.
MiniCPM5 (no native tool template) used a hand-built \texttt{<|im\_start|>} system message;
Qwen3-family models used their native \texttt{<tool\_call>} template. LoRA (r=16, $\alpha$=32,
dropout 0.05) in bf16, merged into standalone checkpoints. The 0.6B recipe was reduced to 1
epoch at lr $5\times10^{-5}$ after a 2-epoch run overfit and produced empty replies.
\begin{table}[h]
\centering
\begin{tabular}{lccc}
\toprule
Model & Version & Tool-call success & Notes\\
\midrule
2B & v2 & 4/4 & direct-answer cases also correct\\
1B & v2 & 4/4 & ---\\
0.6B & v2 & 3/3 & 2-epoch v1 overfit; v2 fixed\\
0.3B & v1 & 3/3 & fixed invented-tool-name failure\\
\bottomrule
\end{tabular}
\caption{Tool-call accuracy on calculator / weather / time.}
\end{table}
A ReAct loop (\texttt{agent\_infer.py}) parses both Qwen-style tool calls and bare nested JSON
via a balanced-brace decoder. Tool prompts require the tool-definition system message.

\section{Related Work and Relationship to JetSpec}
\subsection{Lineage}
\begin{itemize}[leftmargin=*]
  \item Speculative decoding (Leviathan et al., 2023; Chen et al., 2023): cheap draft verified
        by the target with a lossless rule. Amoeba-DA is an instance of this family.
  \item Medusa (Cai et al., 2024): parallel bottleneck draft heads --- direct ancestor of our design.
  \item EAGLE/EAGLE-3 (Li et al., 2024, 2025); DFlash (Chen et al., 2026): better draft alignment.
  \item Early exit: CALM (Schuster et al., 2022), LayerSkip (Elhoushi et al., 2024),
        Mixture-of-Depths; our \S3 is a negative-result study in this tradition.
  \item Cross-layer KV: LCKV (Wu \& Tu, ACL 2024), YOCO (Sun et al., 2024), Cross-Layer
        Attention (Brandon et al., NeurIPS 2024), MLKV (2025) share or compress KV across
        layers, typically requiring pretraining. MiniCache (NeurIPS 2024) merges adjacent KV;
        HShare (ICLR 2025) shares critical-token indices. AttentionPredictor (NeurIPS 2025)
        predicts attention scores, not KV tensors. KV-Direct (arXiv 2603.19664, 2026) proves
        KV are deterministic projections of the residual stream (exact reconstruction) and
        reports failure on sliding-window models -- consistent with our boundary finding.
        Our V4 (\S5) is a frozen-model learned-emulator study with downstream capability
        evaluation, and claims no novelty for KV reconstruction.
\end{itemize}
\noindent\textbf{License note.} LayerSkip code is CC-BY-NC and was not copied; we implemented its
ideas from the paper and used Apache-2.0 CALM as a reference. JetSpec is CC BY 4.0; we cite its
ideas and reuse no code or weights.

\subsection{Independence statement and comparison with JetSpec}
\noindent\textbf{Independent discovery.} The core idea of Amoeba-DA --- parallel draft heads on
a frozen target with batched exact verification --- was developed in our local experiments on
2026-11-22, before we became aware of JetSpec (Hu et al., arXiv:2606.18394v3, submitted
2026-06-25, UCSD / Hao AI Lab). We learned of JetSpec during a later literature survey. We
therefore do \emph{not} claim novelty for the draft-head-and-verify mechanism itself. JetSpec
was published first and reports substantially stronger results; we cite it as the state of the
art and position our work as complementary.

JetSpec trains a causal parallel draft head over fused hidden states from five target layers,
uses a tree-causal mask, constructs a tree under a 256-node budget, and verifies it in one
forward with a custom kernel; it reports up to 9.64$\times$ on MATH-500 with Qwen3-8B on H100,
trained on 780K regenerated samples on 8$\times$H100 (draft cost $c$=0.0085 at budget 16 vs our
$c\approx0.07$).
\begin{table}[h]
\centering
\small
\begin{tabular}{lll}
\toprule
Dimension & Amoeba-DA (this work) & JetSpec (Hu et al., 2026)\\
\midrule
Draft structure & 4 bottleneck MLPs, last hidden & 5-layer causal decoder, fused features\\
Verification & single linear chain, $K$=4 & whole tree, budget 16--256\\
Accepted length $\tau$ & 1.48 (8B) / 2.61 (2B) & up to 10.76\\
Draft cost $c$ & $\sim$0.07 & 0.0085 ($N$=16)\\
Training data & $\sim$30K self-generated tokens, 1$\times$4080S & 780K, 8$\times$H100\\
Objective & hard-label CE & forward-KL distillation\\
Reported 8B speedup & 1.1--1.4$\times$ & 4--9.64$\times$ (H100)\\
\bottomrule
\end{tabular}
\caption{Amoeba-DA versus JetSpec: difference of configuration and budget, not principle.}
\end{table}
\noindent\textbf{Our complementary contributions:} (i) systematic characterisation of why
early-exit dynamic depth fails before adopting verification; (ii) a minimal recipe
(consumer GPU, $\sim$30K tokens, $K$=4) making verification-based acceleration accessible;
(iii) a boundary curve across dense, hybrid-linear and loop architectures; (iv) dynamic loop
routing as an orthogonal dynamic-computation axis.

\section{Negative, Abandoned and Retracted Results}
\begin{enumerate}[leftmargin=*]
  \item Token-level early exit: 16 experiments, all failing quality or speed (\S3).
  \item Math CoT SFT: MATH stayed 0.70 while MMLU fell 0.42$\to$0.20 and HumanEval
        0.6--0.9$\to$0.45. Discarded; no released model contains it.
  \item Draft heads on 0.3B/0.6B: slower (0.36$\times$/0.60$\times$); not shipped.
  \item CCI3 corpus: CE 3.7, $\sim$10$\times$ from deployment distribution; unused in releases.
  \item Qwen3.8-27B-FP8: even after compiling \texttt{causal-conv1d}, $\sim$1 tok/s; abandoned.
  \item Qwen family WITHDRAWN: a training-label offset bug was found and fixed, but end-to-end
        acceptance stayed 0.04--0.07 even though teacher $\alpha_2$ reached 0.76. Crop, replay,
        hidden-match and batch-single checks all failed to restore agreement; the verification
        engine is incompatible with the Qwen KV cache. No Qwen speedup is cited.
  \item Retracted measurement: early 8B run reported 1.62$\times$ with a depressed 8.1 tok/s
        baseline (system load); not cited. All shared-GPU timings are bimodal (2B: 2.8 vs 26
        tok/s; Gemma: 7/23/73 s per prompt; Ling: 1.2 vs 12 tok/s), so reported speedups use
        only single-process stable-fast-mode measurements.
  \item \textbf{vLLM engineering obstacles.} vLLM 0.11 wheels are built against torch 2.8+cu128
        and failed to import under the main torch 2.15+cu132 environment with an undefined-symbol
        error (\texttt{c10\_cuda\_check\_implementation}); downgrading the main environment would
        have broken the compiled \texttt{fla}/\texttt{causal-conv1d} kernels, so we built an isolated
        venv (torch 2.8+cu128) instead. Even then, a transformers-5.x tokenizer incompatibility
        (\texttt{all\_special\_tokens\_extended} missing) required a monkey-patch before
        \texttt{from vllm import LLM}, and several early smoke tests failed only because competing
        GPU jobs left insufficient free VRAM for CUDA-graph capture. In an isolated venv the import,
        model load and generation all work; these were environment issues, not method failures.
  \item The $\sim$1.45$\times$ (45\%) loop result is \emph{dynamic loop routing} (quality-for-speed,
        CE +0.04), NOT Amoeba-DA. Amoeba-DA applied to the loop model gives agreement 1.0
        (415/415) but 0.568$\times$ -- verification too costly. These lines are kept separate.
  \item \textbf{More draft heads (K=8):} extending K from 4 to 8 did not help: holdout
        $\alpha_{6..9}$ fell from 0.195 to 0.117, and expected accepted length dropped from
        2.61 to 1.69 (predicted speedup 2.02$\times\to$1.67$\times$). Draft--target divergence
        grows with offset; adding heads cannot raise acceptance.
  \item \textbf{V2, confidence-layered verification:} splitting drafts by softmax confidence
        (high: accept unverified; mid: single-token verify; low: reject) gave no sweet spot.
        The draft head is badly overconfident (mean $p=0.777$ vs accuracy 0.469): reaching
        95\% precision needs $\tau>0.9999$ covering only 23\% of tokens, so most still need
        verification. At $\tau$ high, 2B reached agreement 1.0 but was \emph{slower than greedy}
        (163\,s vs 117\,s); at $\tau$ low, a single high-confidence error collapsed a sequence.
  \item \textbf{V3 and V4, alternative objectives and layer skips:} forward-KL distillation
        collapsed acceptance ($\alpha_2$: 0.469$\to$0.020) and learned KV prediction preserved
        capability but was slower (0.973$\times$ decode). Both are described in full in \S5,
        together with the K=8 result and the MACs audit.
  \item \textbf{Serving-framework check (positive control):} to test whether speculative
        decoding survives a real serving stack, we ran vLLM 0.11's built-in n-gram speculative
        decoding on MiniCPM5-2B: \textbf{1.806$\times$} (640$\to$1156 tok/s). This confirms the
        mechanism is framework-independent. Our draft heads cannot be plugged into vLLM directly
        because it does not expose hidden states; a Medusa/EAGLE-style proposer adaptation is left
        to future work. This result also shows the absolute throughput gain from the execution
        framework is far larger than any architecture-level acceleration we measured.
  \item \textbf{Gemma numerical-debug dead ends.} Before attributing Gemma's batched-verification
        collapse to numerical non-associativity, we tested and ruled out: attention-mask variants
        (\texttt{ones\_total}, \texttt{ones\_new}, \texttt{none}); explicit position ids; fp32 lm-head
        readout; and changing attention implementation (eager vs sdpa). None changed the per-layer
        hidden divergence (0.125 at layer 0, accumulating to 0.75 by layer 41). Single-token
        verification restored exact agreement (96/96) but was 3.7$\times$ slower. The divergence is
        in the batched matmul itself, not in masking or position encoding.
  \item \textbf{Measurement artifact invalidating speed claims.} Per-prompt baseline timing is
        bimodal and environment-driven, not model-driven: in one run of 20 prompts the first nine
        ran at 2.8 tok/s and the last ten at 26 tok/s, with no difference in prompt content. The
        same artifact appears in capability-evaluation generation (2.5 vs 27 tok/s),
        long-context curves, and shared-GPU runs where baseline and Amoeba independently hit
        fast/slow modes (producing apparent speedups from 0.287$\times$ to 5$\times$). We therefore
        do not report any speedup derived from shared-GPU timing; the few numbers we cite are from
        single-process, warm-up, repeated stable-fast-mode measurements. Several exploratory
        conclusions (e.g.\ early readings that V2/V4 were ``slower'') are weak exactly because of
        this artifact; they are reported as trends, not as precise ratios.
  \item \textbf{Long-context experiment incomplete.} The speedup-vs-context-length curve (128--8K
        tokens) was started but not completed: the measurement process was interrupted and the
        surviving partial data (\texttt{longctx\_phase1.json}) covers only the shortest lengths.
        An early smoke test at 128 tokens showed $\alpha_2\approx0.024$ and $1.16\times$ on two
        prompts -- far below the short-context $\alpha_2=0.469$ -- suggesting substantial decay,
        but we do not present this as a fitted curve. We refer to the short-context limitation in
        \S11 and leave the full curve to future work.
  \item \textbf{Abandoned model targets.} Qwen3.5-4B was downloaded and tested but abandoned:
        baseline 29.9 tok/s, end-to-end 0.887--0.899$\times$, inconsistent batch-forward logits,
        and unusable draft heads. SenseNova-Loop-9B's causal-conv1d kernel would not compile under
        CUDA 12.8, so its baseline was first measured with a reference fallback (5.7 tok/s at loop1,
        2.9 at loop2). SparkX-2.6-4B's released weights are FP8 in most distributions; we located
        a BF16 variant, but its custom code required several transformers-4.57-era compatibility
        patches under transformers 5.24, and its all-FP32 layer execution made baseline slow.
\end{enumerate}

\section{Conclusion}
We reframed ``dynamic activation'' from skipping layers within a token to dynamically
committing several tokens per step after exact verification (Amoeba-DA) and, on loop-structured
models, to dynamically choosing the number of loop iterations. On MiniCPM5-2B the first route
delivers \textbf{2.015$\times$} speedup while preserving downstream capability (79 objective
questions, zero regressions, ROUGE-L 0.94); MiniCPM5-1B gives 1.30$\times$. Our central finding
is that the benefit is bounded by architecture and context: it holds on full-attention dense
models, is numerically fragile on sliding-window hybrids, is blocked by KDA cache semantics on
MoE models, and concentrates in short context because the draft heads are trained on short
on-policy data. All Qwen-family speedups are withdrawn. The orthogonal loop-routing route gives
45\% speedup at 0.04 nats CE cost, a genuine quality-for-speed trade-off distinct from
speculative decoding. We release models (Apache-2.0), draft heads, scripts, raw data and the
full negative-result record -- including withdrawn and failed results -- for reproducibility.

\section{Limitations and Future Work}
\begin{itemize}[leftmargin=*]
  \item \textbf{Short-context training.} All draft heads were trained on short on-policy corpora
        ($\sim$12K tokens). The reported acceleration is measured on short prompts; we do not
        claim it generalizes to long context, and we quantify the decay in \S4.7. Training
        draft heads on long-context on-policy data is future work.
  \item \textbf{Capability scope.} We claim downstream accuracy preservation, not bit-exact
        decoding; token agreement is 0.734 on 2B. Capability is measured on 79 objective QA
        questions; open-ended generation, long-form reasoning, instruction-following nuance and
        safety are not evaluated.
  \item \textbf{Architecture dependence.} The method does not accelerate sliding-window hybrids
        (numerical non-associativity) or KDA/MLA MoE (cache rollback), and Qwen-family numbers
        are withdrawn (\S4.6, \S9). It is not a universal accelerator.
  \item \textbf{Faithfulness--speed trade-off.} Three verification strategies were measured:
        batched verification is fast but numerically non-associative on sliding architectures;
        single-token verification is exact but 0.6--0.7$\times$; confidence layering has no
        sweet spot because draft confidence is miscalibrated. We did not find a scheme that is
        simultaneously faithful and fast, and we state this as an empirical boundary rather
        than an implementation defect (\S9).
  \item \textbf{Layer-skipping ceiling.} A MACs audit and KV-prediction experiments (\S5) show
        that MLP, not attention, is 73.4\% of a layer's cost; $1.5\times$ needs skipping 15--16
        layers, but capability collapses well before that. Learned KV prediction preserves QA
        accuracy yet is slower than the baseline. Alternative layer-emulation designs are future work.
  \item \textbf{Harness dependence.} All reported speedups use HuggingFace \texttt{generate},
        which is Python-bound here ($\sim$100\% CPU, $\sim$3\% GPU): 27 tok/s for 2B versus a
        theoretical bandwidth limit of $\sim$179 tok/s. Under vLLM 0.11 the same model reaches
        64 tok/s single-prompt and $\sim$640--1280 tok/s batched. Our draft heads depend on hidden
        states that vLLM does not expose, so they are bound to the HF harness and their vLLM
        adaptation (e.g. a Medusa/EAGLE-style proposer) is future work. Importantly, the mechanism
        itself is not framework-specific: vLLM's built-in n-gram speculative decoding gives
        \textbf{1.806$\times$} (640$\to$1156 tok/s) on the same model, confirming that
        verification-based multi-token decoding remains effective in a production stack.
        The absolute throughput leap from the execution framework (CUDA graphs, PagedAttention)
        is far larger than any architecture-level gain we measured; we report both facts.
        Shared-GPU timings are bimodal; speedups use single-process stable-fast mode.
  \item Loop routing rests on one model; multi-model validation and token-exit analysis needed.
  \item Combining loop-depth routing with draft verification at the routed depth is future work.
  \item \textbf{Adjacent directions are already occupied (scope note).} Several extensions that
        follow naturally from our negative results have been published by other groups during
        this study; we record them rather than claim them. \emph{From-scratch low-bit training:}
        BitNet b1.58 (Ma et al., arXiv:2402.17764; JMLR 2025) trains ternary weights with a
        straight-through estimator; BitSkip (Bhuvaneswaran \& Liu, arXiv:2510.23766) already
        studies the composition of extreme quantization and early exit, and finds -- consistently
        with our \S3 result on a \emph{pretrained} backbone -- that training from scratch is what
        makes intermediate exits viable. We note that BitNet reports the KV cache of its models
        as losslessly quantizable to 4 bits, but not to 1.58 bits.
        \emph{Low-bit and structured KV caches:} CommVQ (Li et al., ICML 2025) achieves 2-bit
        KV with a RoPE-commutative codebook; VecInfer, AnTKV and KVPath push this to 1--2 bits
        or add a cross-layer dimension. Our ``add one discrete dimension'' idea is an instance
        of this line, not a new mechanism. \emph{Reversible state:} KV-Direct (arXiv:2603.19664)
        reconstructs K/V exactly from one residual per token (5\,KB vs 136\,KB per token) and
        ReplaySSM (NVIDIA, 2026) makes recurrent-state rollback an $O(1)$ buffer-pointer move;
        together these address the verification-failure rollback cost that blocked our
        confidence-layering and KDA experiments (\S9), but they are prior work.
        \emph{Image and multimodal generation is not an unexplored refuge:} early-exit diffusion
        (ASE, Moon et al., ICML 2024; AdaDiff, ECCV 2024; DuoDiff, 2024), dynamic DiT compute
        (DyDiT++, TPAMI 2026; DiffCR, CVPR 2025; MoDDiT, 2026; DC-DiT, 2026), speculative
        diffusion (De Bortoli et al., ICML 2025; SpeCa, 2025; FREE, 2025; SSD, 2026; VVS, 2025)
        and 1--2-bit visual KV caches (CalibQuant, arXiv:2502.14882; AKVQ-VL, 2025; VisCache,
        EMNLP 2026) are all published, often with larger speedups than ours. End-device small
        diffusion models (0.3--0.4B) with elastic depth are likewise studied (SnapGen++, 2026;
        JuZhou 1.0, 2026; EdgeDiT, CVPRW 2026). What remains genuinely underrepresented is not a
        mechanism but an \emph{evaluation discipline}: downstream capability preservation with
        zero regressions instead of bit-exactness or FID alone; boundary mapping across
        architectures that reports failures; and harness-level throughput comparison. TwinPrune
        (Liu et al., arXiv:2610.17530) independently exposed a related evaluation flaw -- VLM
        compression claims validated only on prefill-only probes measure prefill quality, not
        decode quality -- which supports the same view. More broadly, this concern is now an
        active research topic in its own right: ``On Evaluating Performance of LLM Inference
        Systems'' (SoCC 2026) gives a baseline-fairness / setup / metric-design checklist whose
        speculative-decoding case study mirrors our \S4 analysis; ``The Silent Hyperparameter''
        (arXiv:2605.19537) surveys 200 inference engines and 35,000 papers and finds that the
        backend alone shifts benchmark scores by up to 16.6 points, a large-scale confirmation
        of the harness dependence we observed at single-model scale; and a rigor-matched audit
        of layer skipping (arXiv:2608.28846) shows that separating online-search overhead from
        pure inference cost reverses naive wall-clock rankings. These works do not invalidate
        our measurements -- they explain why our evaluation protocol matters. Transferring it to
        diffusion and VLM inference, rather than inventing another accelerator, is the most
        defensible next step.
\end{itemize}

\appendix
\section{Reproducibility}
Key artifacts: \texttt{medusa\_train4.py}, \texttt{medusa\_bench\_final.py},
\texttt{medusa\_agreement.py}, \texttt{agent\_infer.py}, \texttt{gen\_agent\_data2.py},
\texttt{train\_agent\_sft.py}, \texttt{bench\_compare.py},
\texttt{minicpm5\_chat\_template.jinja}, \texttt{ckpt\_\{2b,1b,06b,03b\}\_agent/full\_model/},
\texttt{medusa\_chat.pt}, \texttt{medusa\_chat\_1b.pt}, \texttt{docs/*.md},
\texttt{bench2/final\_results/}.

\section{Default decoding recommendation}
Greedy decoding can loop on long answers (2B repeated a sentence 20+ times). Recommended:
\texttt{temperature=0.8, top\_p=0.95, repetition\_penalty=1.05}. Official model-card scores
use temperature 1.0 and larger budgets, so our short-subset greedy scores are not directly
comparable.

\end{document}
